% Partie : sets-functions-relations
% Chapitre : sets
% Section : important-sets

\documentclass[../../../include/open-logic-section]{subfiles}

\begin{document}

\olfileid[fr]{sfr}{set}{imp}
\olsection{Quelques ensembles importants}

\begin{ex}
Nous étudierons surtout des ensembles dont les !!{element}s sont
des objets mathématiques. Quatre d'entre eux sont assez importants
pour porter un nom particulier :
\begin{multline*}
    \Nat = \{0, 1, 2, 3, \ldots\} \\
    \shoveright{\text{l'ensemble des entiers naturels}}\\
    \shoveleft{\Int = \{\ldots, -2, -1, 0, 1, 2, \ldots\}} \\
    \shoveright{\text{l'ensemble des entiers relatifs}}\\
    \shoveleft{\Rat = \Setabs{\nicefrac{m}{n}}{m, n \in \Int\text{ et }n \neq 0}}\\
    \shoveright{\text{l'ensemble des nombres rationnels}}\\
    \shoveleft{\Real = (-\infty, \infty)}\\
    \text{l'ensemble des nombres réels (le continu)}
\end{multline*}
Ce sont tous des ensembles \emph{infinis} : chacun possède une
infinité d'!!{element}s.

En passant d'un ensemble au suivant, nous disposons de
\emph{nouveaux} nombres. On voit en effet que
$\Nat \subseteq \Int \subseteq \Rat \subseteq \Real$ :
tout entier naturel est un entier relatif ; tout entier relatif
est rationnel ; tout nombre rationnel est réel.
On voit également que $\Nat \subsetneq \Int \subsetneq \Rat$,
puisque $-1$ est un entier relatif qui n'est pas naturel, et
$\nicefrac{1}{2}$ est rationnel sans être entier.
Il est moins évident que $\Rat \subsetneq \Real$, c'est-à-dire
qu'il existe des nombres réels non rationnels\oliflabeldef{sfr:arith:real:realline}{,
mais nous y reviendrons plus loin (\olref[arith][real]{realline})}{}.

Nous utiliserons aussi parfois l'ensemble des entiers strictement
positifs $\PosInt = \{1, 2, 3, \dots\}$ et l'ensemble formé
des deux premiers entiers naturels, $\Bin = \{0, 1\}$.
\end{ex}

\begin{tagblock}{compsci}
\begin{ex}[Mots]
Un autre exemple intéressant est l'ensemble $A^{*}$ des
\emph{mots finis} sur un alphabet $A$ : toute suite finie
d'éléments de~$A$ est un mot sur $A$. Pour tout alphabet~$A$,
le \emph{mot vide $\Lambda$} fait partie des mots sur~$A$.
Par exemple,
\begin{multline*}
\Bin^*
=\{\Lambda,0,1,00,01,10,11,\\
000,001,010,011,100,101,110,111,0000,\ldots\}.
\end{multline*}
Si $x=x_{1}\ldots x_{n}\in A^{*}$ est un mot formé de $n$
«~lettres~» de $A$, sa \emph{longueur} est~$n$,
et on écrit $\len{x}=n$.
\end{ex}
\end{tagblock}

\begin{ex}[Suites infinies]
Pour tout ensemble $A$, on peut aussi considérer l'ensemble~$A^\omega$
des suites infinies d'!!{element}s de~$A$. Une suite infinie
$a_1a_2a_3a_4\dots$ est une liste d'objets qui se prolonge
indéfiniment dans une seule direction, chacun de ces objets
étant !!a{element} de~$A$.
\end{ex}

\end{document}
